\subsection{A genuine two-edge singleton at two outer edges}\label{proof:R2-C047}
\paragraph{R47.1 Original model and exact new statements}

The original target is an orientation of every edge of every finite tree
with at most five vertices of original degree at least three, chosen from
the tree before the host, with original partition at most one for every
finite original probability and every real skew kernel bounded by one.
This section proves the following half-tree estimate, retaining all
original vertices, subdivisions, arms and their repetitions, and the
independent original core lengths.

Work on the unchanged original finite probability space, after deleting
only zero-mass host atoms. Put
\fitdisplay{
Jv=(Ev)1,\quad A=J+K,\quad T=J-K=A^*,\quad
 g=K1,\quad p=1+g,\quad q=1-g,
}{}
\fitdisplay{
f_n=A^{n}\one,\quad \bar f_n=T^{n}\one,\quad h_n=T f_{n-1},\quad
 W_n=A \bar f_{n-1},\qquad \delta_j=\|K^{j}\one\|_2^2,\quad d=\delta_1.
}{}
All expectations, norms, adjoints and spectral measures use this same \(\mu\).
The finite calculus is
\fitdisplay{
\theta=13/32,\quad\kappa=19/32,\quad
 0\le d\le1/3,\quad d^2\le \delta_2\le\theta d,\quad
 \delta_{j+1}\le\theta \delta_j,\quad\|K\|^2\le\theta,
 \quad |Kv(x)|\le\|v\|_2.
}{R47.1.1}
Moment Cauchy also gives \(\delta_2^2\le d \delta_3.\) Define
\fitdisplay{
I_q(U)=E[q^{-1/3}|U|^{8/3}],\qquad
 B_n=E[\bar f_n^4/q],\qquad
 N(d)=(1+d/2)^4(1+32d/19)^2.
}{R47.1.2}
At each retained atom, \(p,q\ge\mu(x)>0\). An actual finite rooted-tree
message is positive, because the all-equal labeling contributes positive
mass. 

Theorem A — exact original weighted two-step transfer. On every such
original host,
\fitdisplay{
\boxed{B_4N(d)\le1-d/100.}
}{R47.1.3}
Consequently every real original function \(X\) satisfies
\fitdisplay{
\boxed{
 I_q(T^2(\bar f_2X))^3
 \le(1-d/100)\bigl(E[\bar f_2|X|^4]\bigr)^2.}
}{R47.1.4}
The weight and operator in (R47.1.4) are part of the statement. It is not a
claim about an ordinary fourth input to a different operator.

Theorem B — every actual \(a=2,\allowbreak  M=2\) singleton. Let an original internal
vertex \(r\) have, after maximal equal-length pairing, exactly one residual
arm of length two. Its original terminal has any strict residual list of
size at least two, its inner connector has any positive length \(L\), and
both vertices may have arbitrarily many genuine equal pairs. Give every
edge the direction in \mbox{\hyperref[part:R47.2]{part R47.2}}. At outer length \(M=2\) the actual message
into the original center obeys
\fitdisplay{
\boxed{I_q(U)^3\le1-d/100\le e^{-d/100}.}
}{R47.1.5}
The same direction table applies on both sides of \(d=6/25\).

Theorem C — exact original assembly. A side from Theorem B can be
combined with any separately supported actual opposite output on the
same \(q\). If that output has \(I_q^3\le1\), the nonreducible-end, nonempty
R-center P5 branch has \(Z\le \exp(-d/800).\) If both sides are from Theorem B,
the corresponding deficit is \(d/400\). With the separate original Q4 and
central-singleton constructions, the latter gives \(Z\le 1\) for arbitrary
original terminal and central multisets, both inner lengths independently
positive, both actual outer lengths two, and both residual internal arms
of length two. No other original length is identified with these twos.

Appendix A gives a separate optional high-domain theorem for \(a\ge 3.\)


\paragraph{R47.2 Complete actual original direction and reference input}\label{part:R47.2}

For an arm, \(f_n\) means every edge outward from its attachment, \(\bar f_n\) every
edge inward, \(h_n\) the first edge inward and the remaining \(n-1\) outward,
and \(W_n\) the first outward and the remaining \(n-1\) inward. These definitions
also specify \(n=1: h_{1}=q\) and \(W_{1}=p.\)

Pair equal arm lengths maximally using a fixed ordering of the original
tree. A genuine pair can be oriented as \(f_l h_l\) in the positive table,
or \(\bar f_l W_l\) in its full reverse. Its factor is in \((0,1]\). The original
two arms are retained, with their actual location and edge count.
The unpaired residual lengths are strictly increasing. 

The complete C034 theorem supplies \(E[f_2(T^L F_{\rm ref})^4]\le1\)
and its genuine-decoration version for all terminal residuals below.
In the table, a canonical pair is \(h_s f_t\) with \(s<t.\) Later pairs are the
remaining consecutive strict residual pairs, each retained once.
Each row gives the outward-reference construction followed by
the actual terminal factor used here after full reversal.

\begin{description}[style=nextline,leftmargin=1.6em]
\item[{Exact \((1,t)\)}] \(F_{\rm ref}=qf_t\); its full reverse is \(p\bar f_t\).
\item[{Exact \((2,t)\)}] \(F_{\rm ref}=h_2f_t\); its full reverse is \(W_2\bar f_t\).
\item[{Exact \((3,t)\), \(L=1\) or \(2\)}] \(F_{\rm ref}=f_3\bar f_t\); its full reverse is \(\bar f_3f_t\).
\item[{Exact \((3,t)\), \(L\ge3\)}] \(F_{\rm ref}=h_3f_t\); its full reverse is \(W_3\bar f_t\).
\item[{Single pair with minimum at least four; every even list with at least four arms}] The product of the canonical pairs \(h_sf_t\); its full reverse is the product of \(W_s\bar f_t\).
\item[{Odd minimum \(\ell\ge2\), second length at least four}] \(\bar f_\ell\) times the later canonical pairs; its full reverse is \(f_\ell\) times the later negative pairs.
\item[{Odd prefix \((2,3,t)\)}] \(\bar f_2f_3h_t\) times the later canonical pairs; its full reverse is \(f_2\bar f_3W_t\) times the later negative pairs.
\item[{Exact \((1,s,t)\), \(s\ge4\)}] \(F_{\rm ref}=qf_sh_t\); its full reverse is \(p\bar f_sW_t\).
\item[{Unit minimum, odd list, second length at least four, at least five arms}] \(q\) times the later canonical pairs; its full reverse is \(p\) times the later negative pairs.
\item[{Odd prefix \((1,3,t)\)}] \(qf_3h_t\) times the later canonical pairs; its full reverse is \(p\bar f_3W_t\) times the later negative pairs.
\item[{Exact \((1,2,3)\), \(L=1\)}] \(F_{\rm ref}=qf_2\bar f_3\); its full reverse is \(p\bar f_2f_3\).
\item[{Exact \((1,2,3)\), \(L\ge2\)}] \(F_{\rm ref}=qf_2h_3\); its full reverse is \(p\bar f_2W_3\).
\item[{Exact \((1,2,t)\), \(t\ge4\)}] \(F_{\rm ref}=qf_2h_t\); its full reverse is \(p\bar f_2W_t\).
\item[{Odd prefix \((1,2,3,4,t_2)\), at least five arms}] \(qf_2h_{t_2}h_3f_4\) times the canonical pairs from index three; its full reverse is \(p\bar f_2W_{t_2}W_3\bar f_4\) times the negative pairs from index three.
\item[{Other odd prefix \((1,2,t)\), at least five arms}] \(qf_2h_t\) times the later canonical pairs; its full reverse is \(p\bar f_2W_t\) times the later negative pairs.
\end{description}

The table is indexed by parity, minimum, second minimum and the
displayed short exceptions. In particular the reference \((1,2,3)\)
exception is at \(L=1.\) It is not the different C106 ordinary-eighth \(L=2\)
exception. A predetermined integer priority resolves overlapping labels.

Let \(E_r^-\) and \(E_t^-\) be the genuine negative pair decorations at \(r\) and
the terminal. Define
\fitdisplay{
X_0=A^L(E_t^-F_{\rm ref}^-),\qquad X=E_r^-X_0,
 \qquad U=T^2(\bar f_2X).
}{R47.2.1}
Thus the sole two-edge arm points toward \(r\), the inner \(L\)-edge connector
points from \(r\) toward its terminal, and both outer edges point from \(r\)
toward the original center. All terminal directions listed as full reversals,
all pair decorations and the genuine reference arm were reversed together.
This is a legal fixed orientation of the complete original half-tree.

The fully reversed C034 reference theorem gives
\fitdisplay{
E[E_r^- \bar f_2 X_0^4]\le1.
}{}
Since \((E_r^-)^4\le E_r^-\) for \(0<E_r^-\le1\),
\fitdisplay{
\boxed{E[\bar f_2X^4]\le1.}
}{R47.2.2}
There is one original \(\bar f_2\) arm in (R47.2.1). The square-root
factorization below uses that arm.

\paragraph{R47.3 Original weighted bridge and its exact norm cost}

For any real \(X\), positivity and row Cauchy for \(T^2M_{\bar f_2}\) give
\fitdisplay{
|T^2(\bar f_2X)|^2\le \bar f_4\,T^2(\bar f_2|X|^2).
}{}
H\"older on the original measure \(\mu\), with exponents \(3\) and \(3/2\),
followed by the ordinary \(L^2\) bound, gives
\fitdisplay{
\begin{aligned}
 I_q(T^2(\bar f_2X))^3
 &\le B_4\{E[T^2(\bar f_2|X|^2)^2]\}^2\\
 &\le B_4\|T^2M_{\sqrt{\bar f_2}}\|^4
                     (E[\bar f_2|X|^4])^2.
 \end{aligned}
}{R47.3.1}
The vector acted on is \(\sqrt{\bar f_2}|X|^2.\) The weight and the square-root
operator are changed together in this exact factorization.

For every positive original operator \(P\) and \(w\ge 0,\allowbreak \) row Cauchy implies
\fitdisplay{
|P(\sqrt w\,v)|^2\le(P|v|)(P(w|v|)).
}{}
After integration and ordinary \(L^2\) Cauchy,
\fitdisplay{
\|PM_{\sqrt w}\|^2\le\|P\|\,\|PM_w\|.
}{R47.3.2}
The actual full-multiplier interface is
\fitdisplay{
\|TM_{\bar f_2}\|^2\le1-d/72.
}{}
Its order follows from full reversal and ordinary adjunction of the
original outward-arm interface, as reconstructed in actual E008 §3.1.
Use \(P=T^2,\allowbreak  w=\bar f_2\) and \(T^2M_{\bar f_2}=T(TM_{\bar f_2}).\) The actual C064 norms give
\fitdisplay{
\begin{aligned}
 \|T^2M_{\sqrt{\bar f_2}}\|^4
 &\le\|T^2\|^2\|T\|^2(1-d/72)\\
 &\le(1+(1024/855)\delta_2)^4(1+32d/19)^2(1-d/72)\\
 &\le N(d).
 \end{aligned}
}{R47.3.3}
The last step uses \(\delta_2\le 13d/32,\allowbreak \) so (\(1024/855)\delta_2\le 416d/855<d/2,\allowbreak \)
and drops only the positive factor \(1-d/72\le 1.\)  Equations (R47.3.1)-(R47.3.3) reduce Theorem A exactly to (R47.1.3).

\paragraph{R47.4 Exact four-edge row and correlated cubic}

Write \(\tau=d-\delta_2,\allowbreak  C=1-d,\allowbreak \) and
\fitdisplay{
R=K G_3(-K)g=Kg-K^2g+K^3g,
 \qquad G_3(z)=1+z+z^2.
}{}
Direct finite multiplication gives
\fitdisplay{
\boxed{\bar f_4=Cq+R-d.}
}{R47.4.1}
Also \(\bar f_4=T^2\bar f_2\le 2(\E\bar f_2)q=2Cq.\) The exact same-vector moments are
\fitdisplay{
ER=-\tau,\qquad h:=E(gR)=\delta_2,\qquad
 V:=ER^2=\delta_2-\delta_3+\delta_4.
}{R47.4.2}
For example \(E(Kg)=-d,\allowbreak  E(K^2g)=0,\allowbreak  E(K^3g)=\delta_2\); the odd cross
inner products vanish by skewness. These are finite identities.

On the original finite spectral block \(z=i \sqrt{x}\),
\(|G_{3}(z)|^2=1-x+x^2\le 1\) for \(0\le x\le \theta.\) Hence
\fitdisplay{
|R|\le\sqrt d,\quad V\le \delta_2,\quad
 |E(gR^2)|\le d\sqrt V\le\tfrac23d^{3/2}.
}{R47.4.3}
The last inequality uses original Cauchy, \(R^2\le d\) and \(\theta<4/9.\)

Set \(W=2R-Cg.\) Its actual precursor is 2G3(-K)g-C1, with
E[G3(-K)g]=d. The original row bound therefore proves
\fitdisplay{
|W|^2\le C^2+4d-4Cd=1-2d+5d^2\le1
 \quad(0\le d\le1/3).
}{R47.4.4}
This improves the looser generic correlated cap specifically at the
actual n4 row. No arbitrary n is relabeled as four.
The genuine mean and second moment are
\fitdisplay{
EW=-2\tau,\qquad EW^2=C^2d+4V-4Ch.
}{}
Since (\(W+2\tau)^2(W-1)\le 0\) pointwise,
\fitdisplay{
EW^3\le(1-4\tau)EW^2-4\tau^2+8\tau^3,
 \qquad E(g^2W)\le d.
}{R47.4.5}

For \(s\ge -1,\allowbreak \)
\fitdisplay{
\min(1+s,2)^4\le1+4s+7s^2+4s^3.
}{R47.4.6}
For -\(1\le s\le 1\) the difference from (\(1+s)^4\) is \(s^2(1-s^2)\ge 0.\)
For \(s\ge 1\) the right side is16 at1 and has positive derivative.
Apply (R47.4.6) at \(s=(R-d)/(Cq),\allowbreak \) which is >-1 by (R47.4.1), keeping
the cap \(\bar f_4\le 2Cq.\) This yields
\fitdisplay{
B_4\le C^4Eq^3+4C^3E[q^2(R-d)]
             +7C^2E[q(R-d)^2]+4CE(R-d)^3.
}{R47.4.7}
The correlated cubic identity used in expanding it is exactly
\fitdisplay{
\begin{aligned}
 &-C^4Eg^3+4C^3E(g^2R)-6C^2E(gR^2)+4CER^3\\
 &\hspace{12mm}=\frac{C}{2}\{EW^3+C^2E(g^2W)\}.
 \end{aligned}
}{R47.4.8}
The coefficient -7 in (R47.4.7) leaves the additional -\(C^2E(gR^2).\)
Substitute (R47.4.2),(R47.4.5) into (R47.4.7)-(R47.4.8). Direct collection gives
\fitdisplay{
\begin{aligned}
 B_4\le{}&K_0+C(9-27d+8\delta_2)(V-h)\\
         &+C(-1+13d-30d^2+8d\delta_2)h-C^2E(gR^2),
 \end{aligned}
}{R47.4.9}
where the complete constant part is
\fitdisplay{
\begin{aligned}
 K_0={}&C^4(1+3d)
 +C^3[(1-2\tau-4d)d-4(d+\tau)]\\
 &+7C^2d^2+14C^2d\tau-12Cd^2\tau-4Cd^3
                     -2C\tau^2+4C\tau^3.
 \end{aligned}
}{R47.4.10}
Every coefficient here comes from the exact n4 identities. In particular
the cap in (R47.4.4) is1, so this is not a transcription of the generic
C100 expression with \(t=1+d/2\) or with beta lowered to kappa d.

\paragraph{R47.5 Complete low-domain row payment, \(0<d\le 6/25\)}

Put \(x=\delta_2/d,\allowbreak \) so \(d\le x\le \theta.\) This is one global moment ratio. Moment Cauchy and (R47.1.1) give
\fitdisplay{
V-h=-\delta_3+\delta_4\le-\kappa \delta_3\le-\kappa x^2d.
}{}
The coefficient 9-27d+8d2 is positive on this domain. Use this bound,
the exact \(h=xd,\allowbreak \) and (R47.4.3) in (R47.4.9). Then
\fitdisplay{
B_4\le C\{1-A(x)d+B(x)d^2-D(x)d^3-E(x)d^4\}
                  +\tfrac23 C^2d^{3/2},
}{R47.5.1}
with
\fitdisplay{
\begin{aligned}
 A(x)&=7-3x+9\kappa x^2,\\
 B(x)&=21-3x+(27\kappa-2)x^2-8\kappa x^3,\\
 D(x)&=20+16x-20x^2+4x^3,\\
 E(x)&=9-2x.
 \end{aligned}
}{R47.5.2}
One direct check of the collection is to write \(v=\tau/d=1-x.\) The
constant part alone satisfies
\fitdisplay{
K_0/C=1-(3+4v)d+(3+20v-2v^2)d^2
                    +(2-26v+4v^3)d^3-(7+2v)d^4.
}{}
Adding the exact h term and the signed V-h term gives precisely (R47.5.2).

The following uniform rational reductions retain enough curvature:
\fitdisplay{
A(x)\ge125/19,\quad B(x)\le22,\quad
 D(x)\ge20+(63/8)d,\quad E(x)\ge8.
}{R47.5.3}
For A, complete its square:
\fitdisplay{
A(x)=9\kappa(x-1/(6\kappa))^2+7-1/(4\kappa)
     =9\kappa(x-1/(6\kappa))^2+125/19.
}{}
For B, its second derivative on [0, theta] is at least \(1055/64>0.\)
Its maximum is therefore at an endpoint; \(B(0)=21\) and
\fitdisplay{
B(\theta)=2854549/131072<22.
}{}
For D use
\(16x-20x^2\ge x(16-20\theta)=63x/8\ge 63d/8\) and \(4x^3\ge 0.\)
Finally \(E\ge 131/16>8.\) Thus (R47.5.1) gives
\fitdisplay{
B_4\le(1-d)[1-(125/19)d+22d^2-20d^3-(127/8)d^4]
                                  +\tfrac23(1-d)^2d^{3/2}.
}{}
Since \(d\le 6/25<1/4,\allowbreak  \sqrt{d}\le 1/2.\) Therefore
\fitdisplay{
\boxed{B_4\le(1-d)R(d),}
 \qquad
 R(d)=1-\frac{356}{57}d+\frac{65}{3}d^2-20d^3-\frac{127}{8}d^4.
}{R47.5.4}

\paragraph{R47.5.1 Exact continuous polynomial certificate}

Let \(P(d)=(1-d)R(d)N(d).\) This is a fixed polynomial. The coefficients of (1-d)R(d), in degrees0 through5, are
\fitdisplay{
(1,-413/57,1591/57,-125/3,33/8,127/8).
}{R47.5.5}
The coefficients of N(d), in degrees0 through6, are
\fitdisplay{
(1,102/19,7995/722,8105/722,34665/5776,588/361,64/361).
}{R47.5.6}
Write \(P=1+\sum_{j=1}^{11}c_j d^j.\) Convolution of (R47.5.5) and (R47.5.6)
gives the following exact coefficients and bounds. Equations (R47.5.5)-(R47.5.6) specify every
coefficient even where a positive expression is more useful than a large
single numerator.


\begingroup\small
\setlength{\tabcolsep}{3pt}
\setlength{\LTleft}{0pt}\setlength{\LTright}{\fill}
\begin{longtable}{@{}p{0.1200\linewidth}p{0.4000\linewidth}p{0.4000\linewidth}@{}}
\caption{A genuine two-edge singleton at two outer edges: parameter values.}\label{tab:parameters:r2-c047:2}\\
\toprule
Entry & Exact coefficient or exact positive description & Bound used \\
\midrule\endfirsthead
\toprule
Entry & Exact coefficient or exact positive description & Bound used \\
\midrule\endhead
\bottomrule\endfoot
1 & \(-107/57\); & retained at the endpoint through P. \\
\addlinespace[2pt]
2 & \(191/2166\); & \(>0\). \\
\addlinespace[2pt]
3 & \(1612016/41154\); & \(>39\). \\
\addlinespace[2pt]
4 & \(4671427/329232\); & \(>0\). \\
\addlinespace[2pt]
5 & \(-50007635/329232\); & \(>-153\). \\
\addlinespace[2pt]
6 & \(-59573162/329232\); & \(>-182\). \\
\addlinespace[2pt]
7 & \(5337151/329232\); & \(>0\). \\
\addlinespace[2pt]
8 & \((127/8)n_3+(33/8)n_4-(125/3)n_5+(1591/57)n_6\); & \(>0\). \\
\addlinespace[2pt]
9 & \((127/8)n_4+(33/8)n_5-(125/3)n_6\); & \(>0\). \\
\addlinespace[2pt]
10 & \((127/8)n_5+(33/8)n_6\); & \(>0\). \\
\addlinespace[2pt]
11 & \((127/8)n_6\); & \(>0\). \\
\end{longtable}
\endgroup


Here \(n_j\) denotes the actual coefficient in (R47.5.6), not a path length.
For c8, the first two terms exceed165+24 and the sole negative term
has magnitude less than70, since \(n_3>11,\allowbreak  n_4>6,\allowbreak  n_5<5/3.\)
For c9 the first term exceeds90 and the negative term has magnitude
less than \(42/5,\allowbreak \) since \(n_6<1/5.\) The last two rows are positive termwise.
The listed integer fractions also show \(c3<40,\allowbreak  c4<15,\allowbreak  c5<-151\) and
\(c6<-180,\allowbreak \) providing further coefficient bounds.

Define \(Q(d)=(P(d)-1)/d\) for \(d>0\) and extend it polynomially at zero.
For every \(0\le d\le 6/25,\allowbreak \) its derivative is bounded below by
\fitdisplay{
\begin{aligned}
 Q'(d)&\ge\frac{191}{2166}+78d-612d^3-910d^4\\
 &\ge\frac{191}{2166}
  +d\left[78-612(6/25)^2-910(6/25)^3\right]\\
 &=\frac{191}{2166}+\frac{471390}{15625}d>0.
 \end{aligned}
}{R47.5.7}
This proves monotonicity on the entire continuous interval. It is not
a statement inferred from checking finitely many d values.

\paragraph{R47.5.2 Exact upper endpoint}

At \(d*=6/25,\allowbreak \) the exact values from (R47.5.4) are
\fitdisplay{
R(d_*)=3116469/7421875,\qquad
 (1-d_*)R(d_*)=3116469/9765625.
}{}
Also
\fitdisplay{
N(d_*)=(28/25)^4(667/475)^2<25/8.
}{R47.5.8}
For an explicit integer check, after multiplying by the positive
denominators (R47.5.8) is
\fitdisplay{
2187629545472<2203369140625.
}{}
Thus
\fitdisplay{
P(d_*)<3116469/3125000=1-8531/3125000.
}{}
By (R47.5.7), \(Q(d)\le Q(d*)<-8531/750000<-1/100.\) Combining with (R47.5.4)
proves
\fitdisplay{
\boxed{B_4N(d)\le1-d/100\qquad(0<d\le6/25).}
}{R47.5.9}
Every expression is for the same actual four-edge path.

\paragraph{R47.6 Complete high-domain row payment, \(6/25\le d\le 1/3\)}

Here a different estimate of that SAME path is more efficient. Its exact
second moment, from (R47.4.1)-(R47.4.2), is
\fitdisplay{
E \bar f_4^2=1-5d+6d^2+d^3+(1-2d)\delta_2-\delta_3+\delta_4.
}{R47.6.1}
For example the cross term in the direct square is
2C E[q(R-d)]=-4Cd, since \(\tau+h=d.\) This retains the true mean.
As in Section5, with \(x=\delta_2/d,\allowbreak \)
\fitdisplay{
(1-2d)x-\kappa x^2\le(1-2d)^2/(4\kappa).
}{}
Consequently
\fitdisplay{
\boxed{E\bar f_4^2\le F(d):=
 1-\frac{87}{19}d+\frac{82}{19}d^2+\frac{51}{19}d^3.}
}{R47.6.2}
No maximizer of this scalar square is imposed on an actual spectral
frequency or on the global moment ratio.

Positivity gives the two simultaneous bounds
\fitdisplay{
b_4\le2(1-d)q,\qquad b_4\le2(1-2d).
}{}
The second uses the exact \(Eb3=1-2d.\) Thus
\fitdisplay{
B_4\le4(1-d)(1-2d)E\bar f_4^2
      \le4(1-d)(1-2d)F(d).
}{R47.6.3}
Let \(H(d)=4(1-d)(1-2d)F(d)N(d).\) The polynomial F is positive
on the entire original domain: its derivative is
\fitdisplay{
F'(d)=(-87+164d+153d^2)/19\le-46/57<0,
 \qquad F(1/3)=1/19.
}{}
In particular \(0<F\le 1\) and F'/\(F\le -46/57.\) On \(d\ge 6/25\) the logarithmic
derivative of H is therefore at most
\fitdisplay{
-25/19-50/13-46/57+2+64/19=-445/741<0.
}{R47.6.4}
All logarithmic factors are positive, so this is a valid continuous
monotonicity proof through \(d=1/3.\)

At \(d*=6/25,\allowbreak \)
\fitdisplay{
F(d_*)=55441/296875,\qquad
 4(1-d_*)(1-2d_*)F(d_*)=2882932/9765625.
}{}
Using the same \(N(d*)<25/8\) from (R47.5.8),
\fitdisplay{
H(d)\le H(d_*)<2882932/3125000<19/20
              <299/300\le1-d/100.
}{R47.6.5}
Together (R47.5.9) and (R47.6.5) prove (R47.1.3) for all \(d>0.\) The estimates apply at their shared endpoint.

At \(d=0,\allowbreak  K1=0\) on the positive support, \(A\one=T\one=\one\) and \(\bar f_4=q=1.\)
Thus (R47.1.3) is equality. The arbitrary-real (R47.1.4) follows directly from
(R47.3.1)-(R47.3.3), also at zero; an arbitrary X need not equal one.
For the actual tree in (R47.2.1), original successive leaf integration does
give \(X=U=q=1,\allowbreak \) even when the balanced skew kernel K is nonzero.
This proves Theorems A and B.
